% Generated draft: requires mathematical and source review.
\documentclass[11pt]{article}
\usepackage[margin=25mm]{geometry}
\usepackage{fontspec}
\setmainfont{Latin Modern Roman}
\usepackage{amsmath,amssymb,mathtools,unicode-math}
\setmathfont{Latin Modern Math}
\usepackage{enumitem,needspace,xurl,hyperref}
\setlength{\parindent}{0pt}
\setlength{\parskip}{5pt}
\setlength{\emergencystretch}{3em}
\newcommand{\N}{\mathbb{N}}
\newcommand{\Z}{\mathbb{Z}}
\newcommand{\Q}{\mathbb{Q}}
\newcommand{\R}{\mathbb{R}}
\newcommand{\C}{\mathbb{C}}
\newcommand{\F}{\mathbb{F}}
\newcommand{\A}{\mathbb{A}}
\newcommand{\PP}{\mathbb{P}}
\newcommand{\E}{\mathbb{E}}
\newcommand{\eps}{\varepsilon}
\newcommand{\dd}{\,\mathrm{d}}
\begin{document}
\section*{UnsolvedMath ID: 30000099}
\textbf{Source problem code:} catalogue entry
\subsection{Definitions and mathematical statement}
An exponential field is a field \(K\) equipped with a homomorphism \(\exp:(K,+)\to(K^\times,\cdot)\). This basic definition alone does not define the pseudo-exponential fields or genericity conditions intended in the source. Let \(\mathcal G\) denote the intended class of generic pseudo-exponential fields with all its source axioms; these axioms are not reproduced in the record.

The first question asks, for each uncountable cardinal \(\lambda\), whether there is a canonical object \(\widetilde K_{\exp}(\lambda)\in\mathcal G\) of cardinality \(\lambda\). If canonical means uniqueness up to exponential-field isomorphism among the intended generic objects, the requested property is
\[
\forall\lambda>\aleph_0\quad
\exists K\in\mathcal G\ (|K|=\lambda)
\quad\text{and}\quad
\forall K,L\in\mathcal G\ (|K|=|L|=\lambda\Rightarrow K\cong L).
\]
This is an interpretation of canonical, not an assertion that every exponential field of that size is isomorphic or that a functorial choice has been constructed. The precise class and meaning of canonical must be recovered from the source.

The second question asks for the stability status of \(\widetilde K_{\exp}\). In a first-order interpretation one would first fix the language and complete theory \(T=\operatorname{Th}(\widetilde K_{\exp})\); \(T\) is \(\kappa\)-stable if, for every parameter set \(A\) of size at most \(\kappa\) in a sufficiently saturated model, the space \(S_1^T(A)\) of complete one-variable types over \(A\) has cardinality at most \(\kappa\). The record does not specify whether stability refers to that first-order theory, an infinitary axiomatization, a reduct, or a different category of structures. Uniqueness in a generic class is not by itself a first-order stability assertion.

The supplied literature triage checked 2026-08-21 says that Zilber reports answers to the canonical-choice and stability questions, but does not reproduce those answers or their hypotheses. It separately describes identification with the complex exponential field as a conjectural consequence. The natural cardinality-compatible identification would concern \(\widetilde K_{\exp}(2^{\aleph_0})\) and \((\C,+,\cdot,\exp)\), not every uncountable size; no theorem or unmentioned hypothesis for this identification is asserted here. That identification is additional context beyond the two questions explicitly extracted.
\subsection{Short English statement}
Ask for a canonical generic pseudo-exponential field in each uncountable cardinality and for its stability classification. The supplied review reports answers to these two questions but leaves the precise axioms and answers unstated; identification with the complex exponential field remains only conjectural in that review.
\subsection{Sources}
\begin{itemize}
\item Local UnsolvedMath registry, permanent ID 30000099, catalogue entry, \emph{Canonical Pseudo-Exponential Fields and Stability}.
\item \emph{Model Theory and Complex Analytic Geometry} (2005), Oberwolfach Reports, pp. 1845--1882, DOI \url{https://doi.org/10.4171/owr/2004/35}. The supplied verification identifies the first two questions in its pseudo-analytic-structures contribution.
\item Zilber, supplied reference label ``exponential sums'': \url{https://doi.org/10.1112/S0024610701002861}.
\item Supplied literature triage dated 2026-08-21, distinguishing reported answers to the two questions from a conjectural complex-exponential identification.
\end{itemize}
\end{document}
